\documentclass[12pt]{article}

% --- Packages ---
\usepackage[margin=1in]{geometry}
\usepackage{amsmath}
\usepackage{amssymb}
\usepackage[T1]{fontenc}
\usepackage[utf8]{inputenc}
\usepackage{hyperref}
\usepackage{titlesec}
\usepackage{enumitem}
\usepackage{booktabs}
\usepackage{microtype}
\usepackage{parskip}
\usepackage{fancyhdr}
\usepackage{array}
\usepackage{longtable}
\usepackage{lscape}
\usepackage{multirow}
\usepackage{caption}
\usepackage{xcolor}
\usepackage{graphicx}

% --- Page style ---
\pagestyle{fancy}
\fancyhf{}
\rhead{\textbf{E3376 $-$ 23}}
\rfoot{\thepage}
\lfoot{\footnotesize Copyright \textcopyright{} ASTM International, 100 Barr Harbor Drive,
PO Box C700, West Conshohocken, PA 19428-2959, United States.}
\renewcommand{\headrulewidth}{0.4pt}
\renewcommand{\footrulewidth}{0.4pt}

% --- Section formatting ---
\titleformat{\section}{\normalfont\bfseries\large}{\thesection.}{0.5em}{}
\titleformat{\subsection}{\normalfont\bfseries}{\thesubsection}{0.5em}{}
\titleformat{\subsubsection}{\normalfont\bfseries}{\thesubsubsection}{0.5em}{}

% --- Annex section numbering ---
\newcounter{annexsec}
\renewcommand{\theannexsec}{A\arabic{annexsec}}

% --- Note environment ---
\newenvironment{astmnote}[1][\relax]{%
  \par\smallskip\noindent\textit{N\textsc{ote}%
  \ifx#1\relax\else~#1\fi---}\ignorespaces
}{%
  \par\smallskip
}

% ============================================================
\begin{document}

% ------ Title block ------
\begin{center}
  {\small Designation: \textbf{E3376 $-$ 23}}\\[1.5ex]
  \rule{\textwidth}{1.5pt}\\[1ex]
  {\LARGE\bfseries Standard Practice for\\[0.4ex]
   Calibration and Usage of Germanium Detectors in\\[0.4ex]
   Radiation Metrology for Reactor Dosimetry$^{1}$}\\[1ex]
  \rule{\textwidth}{1pt}
\end{center}

\noindent
This standard is issued under the fixed designation E3376; the number immediately following
the designation indicates the year of original adoption or, in the case of revision, the year
of last revision. A number in parentheses indicates the year of last reapproval. A superscript
epsilon ($\varepsilon$) indicates an editorial change since the last revision or reapproval.

% ============================================================
\section{Scope}

1.1\quad This standard establishes techniques for calibration, usage, and performance
testing of germanium detectors for the measurement of gamma-ray emission rates of
radionuclides in radiation metrology for reactor dosimetry. The practice is applicable
only to samples of small size, approximating to point sources. It covers the energy and
full-energy peak efficiency calibration as well as the determination of gamma-ray energies
in the 0.06~MeV to 2~MeV energy region and is designed to yield gamma-ray emission rates
with an uncertainty of $\pm 3$\,\% (see Note~1). This technique applies to measurements
that do not involve overlapping peaks, and in which peak-to-continuum considerations are
not important.

\begin{astmnote}[1]
Uncertainty $U$ is given at the 68\,\% confidence level; that is,
\[
  U = \sqrt{\Sigma\sigma_i^2 + \tfrac{1}{3}\Sigma\delta_i^2}
\]
where $\delta_i$ are the estimated maximum systematic uncertainties, and $\sigma_i$ are
the random uncertainties at the 68\,\% confidence level. Other techniques of error
analysis are in use~(1, 2).$^{2}$
\end{astmnote}

1.2\quad Additional information on the setup, calibration, and quality control for
radiometric detectors and measurements is given in IEEE/ANSI N42.14 and in Guide C1402
and Practice D7282.

1.3\quad The values stated in SI units are generally to be regarded as standard. The rad is
an exception.

1.4\quad \textit{This standard does not purport to address all of the safety concerns, if
any, associated with its use. It is the responsibility of the user of this standard to
establish appropriate safety, health, and environmental practices and determine the
applicability of regulatory limitations prior to use.}

1.5\quad \textit{This international standard was developed in accordance with
internationally recognized principles on standardization established in the Decision on
Principles for the Development of International Standards, Guides and Recommendations
issued by the World Trade Organization Technical Barriers to Trade (TBT) Committee.}

% ============================================================
\section{Referenced Documents}

\noindent 2.1\quad \textit{ASTM Standards:}$^{3}$

\begin{itemize}[leftmargin=2.5em,label={}]
  \item C1402\quad Guide for High-Resolution Gamma-Ray Spectrometry of Soil Samples
  \item D7282\quad Practice for Setup, Calibration, and Quality Control of Instruments
        Used for Radioactivity Measurements
  \item E170\quad Terminology Relating to Radiation Measurements and Dosimetry
\end{itemize}

\noindent 2.2\quad \textit{IEEE/ANSI Standard:}$^{4}$

\begin{itemize}[leftmargin=2.5em,label={}]
  \item N42.14\quad Calibration and Usage of Germanium Detectors Spectrometers for
        Measurement of Gamma-Ray Emission Rates of Radionuclides
\end{itemize}

% ============================================================
\section{Terminology}

3.1\quad \textit{Definitions:}

3.1.1\quad \textit{certified radioactivity standard source}---a calibrated radioactive
source, with stated uncertainties, whose calibration is certified by the source supplier
as traceable to an international or national standards laboratory.

3.1.2\quad \textit{check source}---a radioactivity source, not necessarily calibrated,
that is used to confirm the continuing satisfactory operation of an instrument.

3.1.3\quad \textit{correlated photon summing}---the simultaneous detection of two or more
photons originating from a single nuclear disintegration.

3.1.4\quad \textit{dead time}---the time after a triggering pulse during which the system
is unable to retrigger.

3.1.5\quad \textit{FWHM}---(full width at half maximum) the full width of a gamma-ray
peak distribution measured at half the maximum ordinate above the continuum.

3.1.6\quad \textit{FWTM}---(full width at tenth maximum) the full width of a gamma-ray
peak distribution measured at one tenth of the maximum ordinate above the continuum.

3.1.7\quad \textit{national radioactivity standard source}---a calibrated radioactive
source prepared and distributed as a standard reference material by an international or
national standards laboratory.

3.1.7.1\quad \textit{Discussion}---In the United States, the national standards laboratory
is the U.S. National Institute of Standards and Technology (NIST).

3.1.8\quad \textit{resolution, gamma ray}---the measured FWHM, after background
subtraction, of a gamma-ray peak distribution, expressed in units of energy.

3.2\quad \textit{Acronyms:}

3.2.1\quad \textit{GUM}---the GUM, Guide to the Expression of Uncertainty in
Measurement~(2).

3.2.2\quad \textit{ROI}---region of interest.

3.3\quad For other relevant terms, see Terminology E170.

\begin{astmnote}[2]
The terms ``standard source'' and ``radioactivity standard'' are general terms used to
refer to the sources and standards of National Radioactivity Standard Source and Certified
Radioactivity Standard Source.
\end{astmnote}

% ============================================================
\section{Summary of Practice}

4.1\quad This practice describes the calibration and usage of germanium detectors for
measurement of gamma-ray emission rates of radionuclides in radiation metrology for
nuclear reactors. The practice is intended for use by knowledgeable persons who are
responsible for the development of correct procedures for the calibration and usage of
germanium detectors.

4.2\quad A source emission rate for a gamma ray of a selected energy is determined from
the counting rate in a full-energy peak of a spectrum, together with the measured
efficiency of the spectrometry system for that energy and source location. It is usually
not possible to measure the efficiency directly with emission rate standards at all desired
energies. Therefore, a curve or function is constructed to permit interpolation between
available calibration points.

% ============================================================
\section{Significance and Use}

5.1\quad High-purity germanium detectors are used for precise gamma-ray spectroscopy for
the purpose of determining radioactivity in materials. Typical applications include
monitoring, mapping, and characterization of neutron energy spectra in nuclear reactors or
isotopic fission sources.

% ============================================================
\section{Apparatus}

6.1\quad A typical gamma-ray spectrometry system used in reactor dosimetry consists of a
germanium crystal, cooling system (cryogenic or electronic), power supply, and either
digital or analog signal processing. When gamma rays interact with the detector, they
produce pulses which must be analyzed and counted. Digital processing is often integrated
into the detector system with fast analog-to-digital converter, mathematical shaping,
signal, and amplification followed by the binning of the pulses. In analog systems, signal
processing takes place in individual modules. In both cases, the resultant channelized
data are collected and stored electronically for further analysis.

% ============================================================
\section{Preparation of Apparatus}

7.1\quad Follow the manufacturer's instructions for setting up and preliminary testing of
the equipment. Observe all of the manufacturer's limitations and cautions. All tests
described in Section~12 should be performed before starting the calibrations, and all
corrections made when required. A check source may be used to check the stability of the
system at least before and after the calibration. In cases where electronic settings are
adjustable by the user, it is essential that all the settings be the same for calibration
and subsequent use.

% ============================================================
\section{Calibration Procedure}

8.1\quad \textit{Energy Calibration}---Determine the energy calibration (channel number
versus gamma-ray energy) of the detector system at a fixed gain by determining the
channel numbers corresponding to full-energy peak centroids from gamma rays emitted over
the full-energy range of interest from multipeaked or multinuclide radioactivity sources,
or both. Determine nonlinearity correction factors as necessary~(3).

8.1.1\quad Using suitable gamma-ray compilations~(4, 5), plot or fit to an appropriate
mathematical function the values for peak centroid (in channels) versus gamma energy.

8.2\quad \textit{Efficiency Calibration:}

8.2.1\quad Accumulate an energy spectrum using calibrated radioactivity standards at a
desired and reproducible source-to-detector distance. Set the ROI for each peak of
interest to include at least all channels within the FWTM of the peak. At least 20\,000
net counts should be accumulated in each full-energy gamma-ray peak of interest using
National or Certified Radioactivity Standard Sources, or both (see 12.1, 12.5, and 12.6).

8.2.2\quad For each standard source, obtain the net count rate (total count rate of region
of interest minus the Compton continuum count rate and, if applicable, the ambient
background count rate within the same region) in the full-energy gamma-ray peak, or
peaks, using a method that provides consistent results (see 12.2, 12.3, and 12.4).

8.2.3\quad Correct the standard source emission rate for decay to the date and time of the
count in 8.2.2.

8.2.4\quad Calculate the full-energy peak efficiency, $E_f$, as follows:

\begin{equation}
  E_f = \frac{N_p}{N_\gamma}
  \label{eq:efficiency}
\end{equation}

\noindent where:
\begin{align*}
  E_f      &= \text{full-energy peak efficiency (counts per gamma ray emitted),}\\
  N_p      &= \text{net gamma-ray count in the full-energy peak (counts per second}\\
           &\quad\text{live time) (Note 3) (see 8.2.2), and}\\
  N_\gamma &= \text{gamma-ray emission rate (gamma rays per second).}
\end{align*}

\begin{astmnote}[3]
Any other unit of time is acceptable provided it is used consistently throughout.
\end{astmnote}

8.2.5\quad There are many ways of calculating the net gamma-ray count. This practice
specifies a valid, common technique, applicable when there are no interferences from photo
peaks within an energy range at least equal to the FWTM adjacent to the peak of interest.

8.2.5.1\quad Other net peak area calculation techniques used for single peaks when there
is interference from adjacent peaks are not considered in this practice. Other techniques
are acceptable, if they are used in a consistent manner and have been verified to provide
accurate results.

8.2.5.2\quad Using a simple model, the net peak area for a single peak can be calculated
as follows:

\begin{equation}
  N_A = G_s - B - I
  \label{eq:net_peak_area}
\end{equation}

\noindent where:
\begin{align*}
  G_s &= \text{gross count in the peak region of interest (ROI) in the sample spectrum,}\\
  B   &= \text{continuum, and}\\
  I   &= \text{number of background counts in the region of interest}\\
      &\quad\text{(if there is no background peak, or if a background}\\
      &\quad\text{subtraction is not performed, }I=0\text{).}
\end{align*}

8.2.5.3\quad The net gamma-ray count, $N_p$, is related to the net peak area as follows:

\begin{equation}
  N_p = N_A / T_s
  \label{eq:net_count_rate}
\end{equation}

\noindent where $T_s$ is the spectrum live time.

8.2.5.4\quad The continuum, $B$, is calculated from the sample spectrum using the
following equation (see Fig.~\ref{fig:spectral_peak}):

\begin{equation}
  B = \frac{N}{2n}\!\left(B_{1s} + B_{2s}\right)
  \label{eq:continuum}
\end{equation}

\noindent where:
\begin{align*}
  N     &= \text{number of channels in the peak ROI,}\\
  n     &= \text{number of continuum channels on each side,}^{5}\\
  B_{1s}&= \text{sum of counts in the low-energy continuum region in the sample spectrum, and}\\
  B_{2s}&= \text{sum of counts in the high-energy continuum region in the sample spectrum.}
\end{align*}

\begin{astmnote}[4]
These equations assume that the channels that are used to calculate the continuum do not
overlap with the peak ROI, and are adjacent to it or have the same size gap between the
two regions on both sides. A different equation must be used if the gaps are of a
different size.
\end{astmnote}

8.2.5.5\quad The peak background, $I$, is calculated from a separate background
measurement using the following equation:

\begin{equation}
  I = \frac{T_s}{T_b}\,I_b
  \label{eq:peak_background}
\end{equation}

\noindent where:
\begin{align*}
  T_s &= \text{live time of the sample spectrum,}\\
  T_b &= \text{live time of the background spectrum, and}\\
  I_b &= \text{net background peak area in the background spectrum.}
\end{align*}

8.2.5.6\quad If a separate background measurement exists, the net background peak area is
calculated from the following equation:

\begin{equation}
  I_b = G_b - B_b
  \label{eq:net_bg_area}
\end{equation}

\noindent where:
\begin{align*}
  G_b &= \text{sum of gross counts in the background peak region (of the background spectrum), and}\\
  B_b &= \text{continuum counts in the background peak region (of the background spectrum).}
\end{align*}

8.2.5.7\quad The continuum counts in the background spectrum are calculated from the
following equation:

\begin{equation}
  B_b = \frac{N}{2n}\!\left(B_{1b} + B_{2b}\right)
  \label{eq:bg_continuum}
\end{equation}

\noindent where:
\begin{align*}
  N     &= \text{number of channels in the background peak ROI,}\\
  n     &= \text{number of continuum channels on each side (assumed to be the same on both sides),}\\
  B_{1b}&= \text{sum of counts in the low-energy continuum region in the background spectrum, and}\\
  B_{2b}&= \text{sum of counts in the high-energy continuum region in the background spectrum.}
\end{align*}

8.2.5.8\quad If the standard source is calibrated in units of Becquerels, the gamma-ray
emission rate is given by:

\begin{equation}
  N_\gamma = A\,P_\gamma
  \label{eq:emission_rate}
\end{equation}

\noindent where:
\begin{align*}
  A       &= \text{number of nuclear decays per second, and}\\
  P_\gamma &= \text{probability per nuclear decay for the gamma ray~(4, 5).}
\end{align*}

8.2.6\quad In cases where a calibrated standard source is available for a particular
isotope that is also used in an unknown sample, the measured efficiency for that isotope
is used directly for the unknown.

8.2.7\quad Plot, or fit to an appropriate mathematical function, the values for
full-energy peak efficiency (determined in 8.2.4) versus gamma-ray energy (see 12.5)
(6--14).

\begin{figure}[ht]
\centering
\fbox{\rule{0pt}{2.5in}\rule{4in}{0pt}}
\caption{Typical Spectral Peak with Parameters Used in the Peak Area Determination
Indicated. The figure shows a peak with gross area $G_S$, net area $N_A$, continuum
$B$, and flanking continuum regions $B_1$ and $B_2$ each spanning $n$ channels,
with the ROI spanning $N$ channels.}
\label{fig:spectral_peak}
\end{figure}

% ============================================================
\section{Measurement of Gamma-Ray Emission Rate of the Sample}

9.1\quad Place the sample to be measured at the source-to-detector distance used for
efficiency calibration (see 12.6).

9.1.1\quad Accumulate the gamma-ray spectrum, recording the count duration.

9.1.2\quad Determine the energy of the gamma rays present by use of the energy
calibration obtained under and at the same gain as 8.1.

9.1.3\quad Obtain the net count rate in each full-energy gamma-ray peak of interest as
described in 8.2.2 using the same region of interest and continuum regions as were used
in the efficiency measurement. In cases where the width of the region of interest varies
with energy, use a consistent method for choosing the regions.

9.1.4\quad Determine the full-energy peak efficiency for each energy of interest from the
curve or function obtained in 8.2.5.

9.1.5\quad Calculate the number of gamma rays emitted per unit live time for each
full-energy peak as follows:

\begin{equation}
  N_\gamma = \frac{N_p}{E_f}
  \label{eq:emission_from_count}
\end{equation}

9.1.5.1\quad When calculating a nuclear transmutation rate from a gamma-ray emission
rate determined for a specific radionuclide, a knowledge of the gamma-ray probability per
decay is required~(4, 5), that is,

\begin{equation}
  A = \frac{N_\gamma}{P_\gamma}
  \label{eq:transmutation_rate}
\end{equation}

9.1.6\quad Calculate the net peak area statistical uncertainty as follows:

\begin{equation}
  S_{N_A} = \sqrt{G_s + \left(\frac{N}{2n}\right)^{\!2}\!\left(B_{1s}+B_{2s}\right)
    + \left(\frac{T_s}{T_b}\right)^{\!2}\!\left(S_{I_b}\right)^2}
  \label{eq:stat_uncertainty}
\end{equation}

\noindent where:

\begin{equation}
  S_{I_b} = \sqrt{G_b + \left(\frac{N}{2n}\right)^{\!2}\!\left(B_{1b}+B_{2b}\right)}
  \label{eq:bg_uncertainty}
\end{equation}

\noindent where:
\begin{align*}
  S_{N_A} &= \text{net peak area statistical uncertainty (at }1\sigma\text{ confidence level),}\\
  G_s     &= \text{gross counts in the peak ROI of the sample spectrum,}\\
  G_b     &= \text{gross counts in the peak ROI of the background spectrum,}\\
  N       &= \text{number of channels in the peak ROI,}\\
  n       &= \text{number of continuum channels on each side (assumed to be the same}\\
           &\quad\text{on both sides for these equations to be valid),}\\
  B_{1s}  &= \text{continuum counts left of the peak ROI in the sample spectrum,}\\
  B_{2s}  &= \text{continuum counts right of the peak ROI in the sample spectrum,}\\
  B_{1b}  &= \text{continuum counts left of the peak ROI in the background spectrum,}\\
  B_{2b}  &= \text{continuum counts right of the peak ROI in the background spectrum,}\\
  T_s     &= \text{live time of the sample spectrum, and}\\
  T_b     &= \text{live time of the background spectrum.}
\end{align*}

9.1.6.1\quad If there is no separate background measurement, or if no background subtract
is performed, $S_{I_b} = 0$.

9.1.7\quad For other sources of uncertainty, see Section~11.

% ============================================================
\section{Performance Testing}

10.1\quad System tests should be performed on a regularly scheduled basis (or, if
infrequently used, preceding the use of the system). The frequency for performing each
test will depend on the stability of the particular system as well as on the accuracy and
reliability of the required results. Where health or safety is involved, more frequent
checking may be appropriate. A range of typical frequencies for noncritical applications
is given below for each test.

10.1.1\quad Check the system energy calibration (typically daily to weekly), using two or
more gamma rays whose energies span at least 50\,\% of the calibration range of interest.
Correct the energy calibration, if necessary. Sample counting must be halted or redone if
the system energy calibration is found to be inadequate.

10.1.2\quad Check the system count rate reproducibility (typically weekly to monthly)
using at least one long-lived radionuclide. Correct for radioactive decay if significant
decay ($>1$\,\%) has occurred between checks.

10.1.3\quad Check the system resolution (typically weekly to monthly) using at least one
gamma-ray emitting radionuclide.

10.1.4\quad Check the efficiency calibration (typically monthly to yearly) using a
National or Certified Radioactivity Standard (or Standards) emitting gamma rays with
energies spanning the intended range of use.

10.2\quad Results of all performance checks should be recorded in such a way that
deviations from the norm will be readily observable. Appropriate action, which could
include confirmation, repair, and recalibration as required, should be taken when the
measured values fall outside the predetermined limits.

10.2.1\quad In addition, the above performance checks (see 10.1) are recommended after an
event (such as power failures or repairs) which might lead to potential changes in the
system.

% ============================================================
\section{Sources of Uncertainty}

11.1\quad Other than Poisson distribution uncertainties, the principal sources of
uncertainty (and typical magnitudes) in this technique are:

11.1.1\quad The calibration of the standard source, including uncertainties introduced in
using a standard radioactivity solution, or aliquot thereof, to prepare another (working)
standard for counting, typically $\pm 1$ to $2$\,\% (1S\,\%).

11.1.2\quad The reproducibility in the determination of net full-energy peak counts,
typically $\pm 2$\,\% (1S\,\%).

11.1.3\quad The reproducibility of the positioning of the source relative to the detector
and the source geometry, typically $\pm 1$\,\% (1S\,\%).

11.1.4\quad The accuracy with which the full-energy peak efficiency at a given energy can
be determined from the calibration curve or function, typically $\pm 1.5$\,\% (1S\,\%).

11.1.5\quad The accuracy of the live-time determinations and pile-up corrections,
typically $\pm 1$\,\% (1S\,\%).

% ============================================================
\section{Precautions and Tests}

12.1\quad \textit{Random Summing and Dead Time:}

12.1.1\quad \textit{Precaution}---The shape and length of pulses used can cause a
reduction in peak areas due to random summing of pulses at rates of over a few hundred
per second~(15, 16). Sample count rates should be low enough to reduce the effect of
random summing of gamma rays to a level where it may be neglected, or one should use
pile-up rejectors and live-time circuits, or reference pulser techniques of verified
accuracy at the required rates~(17--23).

\begin{astmnote}[5]
Use of percent dead time to indicate whether random summing can be neglected may not be
appropriate.
\end{astmnote}

12.1.2\quad \textit{Test:}

12.1.2.1\quad Accumulate a $^{60}$Co spectrum with a total count rate of less than 1000
counts per second until at least 25\,000 counts are collected in the 1.332~MeV and
1.173~MeV full-energy peaks. A mixed isotopic point source may be used. Record the
counting live time. The source may be placed at any convenient distance from the
detector.

12.1.2.2\quad Evaluate the activity of $^{60}$Co utilizing first the full photon peak
area at 1.332~MeV and then the area at 1.173~MeV, including any techniques employed to
correct for pile-up and dead time losses.

12.1.2.3\quad Without moving the $^{60}$Co source, introduce a second source, such as
$^{57}$Co, having no gamma rays emitted with an energy greater than 0.662~MeV. Position
the added source so that the highest count rate used for gamma-ray emission rate
determinations has been achieved.

12.1.2.4\quad Erase the first spectrum and accumulate another spectrum for the same
length of time as in 12.1.2.1. The same live time may be used, if the use of live time
constitutes at least a part of the correction technique.

12.1.2.5\quad Evaluate the activity of $^{60}$Co utilizing first the full photon peak
area at 1.332~MeV and then the area at 1.173~MeV, including any techniques employed to
correct for pile-up and dead time losses. For the correction technique to be acceptable,
the resolution must not have increased beyond the range of the technique and the
corrected activity will differ from those in 12.1.2.2 by no more than 2\,\% $1\sigma$
(67\,\% confidence level).

12.2\quad \textit{Peak Evaluation:}

12.2.1\quad \textit{Precaution}---Many techniques~(24--29) exist for specifying the
full-energy peak area and removing the contribution of any continuum under the peak.
Within the scope of this standard, various techniques give equivalent results if they are
applied consistently to the calibration standards and the sources to be measured, and if
they are not sensitive to moderate amounts of underlying continuum. A test of the latter
point is a recommended part of this technique.

12.2.2\quad \textit{Test:}

12.2.2.1\quad Accumulate a spectrum from a mixed isotopic point source until at least
20\,000 net counts are recorded in the peaks of interest lower in energy than 0.662~MeV.
The source may be placed at any convenient distance from the detector.

12.2.2.2\quad Determine the net peak areas of the peaks chosen in 12.2.2.1 with the
technique to be tested. Include any calculations employed by the analysis technique to be
tested to correct for dead time losses, pile-up, and background contributions.

12.2.2.3\quad Without moving the mixed isotopic point source, introduce a $^{137}$Cs,
$^{60}$Co, or any other source with no full-energy photons emitted with energies in the
range 0.060~MeV to 0.600~MeV so the continuum level of the spectrum in this range is
increased 20 times.

12.2.2.4\quad Erase the first spectrum and accumulate another spectrum for the same live
time as in 12.2.2.1, if the use of live time constitutes at least a part of the
correction technique.

12.2.2.5\quad Determine the net peak areas of the same peaks chosen in 12.2.2.1 with the
technique to be tested. Include any calculations employed by the analysis technique to be
tested to correct for dead time losses, pile-up, and background contributions.

12.2.2.6\quad Ideally, the deviations of the 12.2.2.5 net peak areas from the 12.2.2.2
values will be no more than 2\,\% $1\sigma$ (67\,\% confidence level) for the evaluation
technique to be acceptable.

12.3\quad \textit{Correlated Photon Summing Correction:}

12.3.1\quad Correlated photon summing corrections are not needed when an instrument is
calibrated using a standard source of the same radioisotope as the unknown sample.

12.3.2\quad If an energy-dependent efficiency calibration is performed at a large enough
source-detector distance, correlated photon summing corrections may not be needed for
sources that produce multiple photons in cascade. In such a case, isotope-specific
calibrations can be obtained at closer distances by a ratio method using a check source
of each isotope at large and small distance.

12.3.3\quad When another gamma ray or X-ray is emitted in cascade with the gamma ray
being measured, in many cases a multiplicative correlated summing correction, $C$, must
be applied to the net full-energy-peak count rate if the sample-to-detector distance less
than 1.5 times the diameter of the detector. The correction factor is expressed as:

\begin{equation}
  C = \frac{1}{\prod_{i}^{n}(1 - q_i\,\varepsilon_i)}
  \label{eq:summing_correction}
\end{equation}

\noindent where:
\begin{align*}
  C          &= \text{correlated summing correction to be applied to the measured count rate,}\\
  n          &= \text{number of gamma or X-rays in correlation with gamma ray of interest,}\\
  i          &= \text{identification of correlated photon,}\\
  q_i        &= \text{fraction of the gamma ray of interest in correlation with the }i\text{th photon, and}\\
  \varepsilon_i &= \text{total detection efficiency of }i\text{th correlated photon.}
\end{align*}

12.3.3.1\quad Correlated summing correction factors for the primary gamma rays of
radionuclides $^{60}$Co, $^{88}$Y, and $^{46}$Sc are approximately 1.09 and 1.03 for a
65~cm$^3$ detector at 1~cm and at 4~cm sample-to-detector distances, respectively, and
approximately 1.01 for a 100~cm$^3$ detector at a 10~cm sample-to-detector distance. The
$q_i$ must be obtained from the nuclear decay scheme, while the $\varepsilon_i$, which
are slowly varying functions of the energy, can be measured or calculated~(30--32).

12.3.4\quad A similar correction must be applied when a weak gamma ray occurs in a decay
scheme as an alternate decay mode to two strong cascade gamma rays with energies that
total to that of the weak gamma ray~(33). The correction is over 5\,\% for the 0.40~MeV
gamma ray of $^{75}$Se when a source is counted 10~cm from a 65~cm$^3$ detector. Other
common radionuclides with similar-type decay schemes, however, do not require a
correction of this magnitude. For example, $^{47}$Ca (1.297~MeV), $^{59}$Fe (1.292~MeV),
$^{144}$Pr (2.186~MeV), $^{187}$W (0.686~MeV), and $^{175}$Yb (0.396~MeV) require
corrections between 0.990 and 0.998 when counted at 4~cm from a 65~cm$^3$ detector.

12.4\quad \textit{Correction for Decay During the Counting Period:}

12.4.1\quad If the value of a full-energy peak counting rate is determined by a
measurement that spans a significant fraction of a half-life, and the value is assigned
to the beginning of the counting period, a multiplicative correction, $F_b$, must be
applied:

\begin{equation}
  F_b = \frac{\lambda t}{1 - e^{-\lambda t}}
  \label{eq:decay_correction_begin}
\end{equation}

\noindent where:
\begin{align*}
  F_b      &= \text{decay during count correction (count rate referenced to beginning of}\\
           &\quad\text{counting period),}\\
  t        &= \text{elapsed counting time,}\\
  \lambda  &= \text{radionuclide decay constant }\!\left(\dfrac{\ln 2}{T_{1/2}}\right)\text{, and}\\
  T_{1/2}  &= \text{radionuclide half-life.}
\end{align*}

$t$ and $T_{1/2}$ must be in the same units of time ($F_b = 1.01$ for $t/T_{1/2} = 0.03$).

12.4.2\quad If under the same conditions the counting rate is assigned to the midpoint of
the counting period, the multiplicative correction $F_m$ will be essentially 1 for
$t/T_{1/2} = 0.03$ and 0.995 for $t/T_{1/2} = 0.5$. If it need be applied, the
correction to be used is:

\begin{equation}
  F_m = \frac{\lambda t}{1 - e^{-\lambda t}}\,e^{-\lambda t/2}
  \label{eq:decay_correction_mid}
\end{equation}

12.5\quad \textit{Efficiency versus Energy Function or Curve}---The expression or curve
showing the variation of efficiency with energy (see Fig.~\ref{fig:efficiency_curve} for
an example) must be determined for a particular detector~(6--14), and must be checked for
changes with time as specified in the standard. If the full energy range covered by this
standard is to be used, calibrations should be made at least every 0.1~MeV from 0.06~MeV
to 0.30~MeV, about every 0.2~MeV from 0.3~MeV to 1.4~MeV, and at least at one energy
between 1.4~MeV and 2~MeV. Radionuclides emitting two or more gamma rays with
well-established relative gamma-ray probabilities may be used to better define the form
of the calibration curve or function. A calibration with the same radionuclides that are
to be measured should be made whenever possible and may provide the only reliable
calibration when a radionuclide with cascade gamma rays is measured very close to the
detector.

12.6\quad See mandatory Annex~A1 for the fitting method to be used for determination of
the energy-dependent function to be used for detector efficiency and for the estimation
of the correlation coefficient between efficiencies at pairs of energies.

12.7\quad \textit{Source Geometry}---A gamma ray undergoing even small-angle scattering
is lost from the narrow full-energy peak, making the full-energy peak efficiency sensitive
to the source or container thickness and composition. For most accurate results, the
source to be measured must duplicate as closely as possible the calibration standards in
all aspects (for example, shape, physical and chemical characteristics, etc.). If this is
not practicable, appropriate corrections must be determined and applied.

12.7.1\quad If the source shape and detector distance remain constant, attenuation within
the source material is corrected as follows:

\begin{equation}
  A_c = A_o\,\frac{\mu x}{1 - e^{-\mu x}}
  \label{eq:source_attenuation}
\end{equation}

\noindent where:
\begin{align*}
  A_c &= \text{corrected number of nuclear decays per second,}\\
  A_o &= \text{observed number of nuclear decays per second,}\\
  \mu &= \text{cm}^2/\text{g} = \text{mass attenuation coefficient~(34), and}\\
  x   &= \text{g}\cdot\tfrac{1}{\text{cm}^3}\cdot\text{cm}
        = \text{mass times path length divided by volume.}
\end{align*}

\begin{figure}[ht]
\centering
\fbox{\rule{0pt}{3in}\rule{4.5in}{0pt}}
\caption{Typical Efficiency versus Energy for a Germanium Detector. The vertical
axis shows absolute full-energy peak efficiency (log scale, 0.01 to 1) and the
horizontal axis shows energy in MeV (log scale, 0.050 to 2.000). The curve rises
to a maximum near 0.1--0.2~MeV and falls smoothly at higher energies.}
\label{fig:efficiency_curve}
\end{figure}

12.7.2\quad If the source shape, composition, and detector distance remain constant, the
attenuation of an interposed absorber is corrected as follows:

\begin{equation}
  A_c = A_o\cdot e^{\mu x}
  \label{eq:absorber_attenuation}
\end{equation}

% ============================================================
\section{Precision and Bias}

\begin{astmnote}[6]
Measurement uncertainty is described by a precision and bias statement in this standard.
Another acceptable approach is to use Type~A and Type~B uncertainty components~(1, 2).
The Type~A/B uncertainty specification is now used in the International Organization for
Standardization (ISO) standards and this approach can be expected to play a more
prominent role in future uncertainty analyses.
\end{astmnote}

13.1\quad Tests in individual laboratories have indicated that activities can be determined
with a bias of $\pm 3$\,\% (1S\,\%) and with a precision of $\pm 1$\,\% (1S\,\%).

% ============================================================
\section{Keywords}

14.1\quad calibration of germanium detectors; germanium detectors; radiation metrology;
reactor dosimetry; usage of germanium detectors

% ============================================================
% ANNEX
% ============================================================
\newpage
\begin{center}
  \rule{\textwidth}{1pt}\\[0.5ex]
  {\large\bfseries ANNEX}\\[0.3ex]
  {\normalsize\bfseries (Mandatory Information)}\\[0.5ex]
  \rule{\textwidth}{1pt}
\end{center}

\section*{A1.\quad FITTING OF ENERGY-DEPENDENT FORMULAE FOR FULL-ENERGY
EFFICIENCY OF HPGe DETECTORS}

\begin{astmnote}[A1.1]
Material in this annex is based on Ref~(35). General guidance on the linear model for
least-squares fitting can be found in Ref~(36).
\end{astmnote}

\subsection*{A1.1\quad Overview}

A1.1.1\quad A calibration of the full-energy detection efficiency of gamma spectrometers
is performed using multi-element standard sources, or a combination of these together
with sources of some isotopes that produce many gamma lines. The sources used do not
include all isotopes which the detectors are later used to measure, and therefore the
calibration has to include an energy-dependent fitting to obtain the efficiencies at
energies other than those in the calibration measurements.

A1.1.2\quad This fitting process introduces strong correlations between gamma lines of
nearby energies.

A1.1.3\quad Examples are given in this annex of detector calibrations in which the
correlations are calculated and propagated through an analysis of measured activities.
Methods of calculating the correlations are described in detail, and numerical examples
are presented in which strong correlations are present between activities measured for
different activation detectors.

\begin{astmnote}[A1.2]
Numerical values given in this annex are for illustrative purposes only.
\end{astmnote}

A1.1.4\quad The effect of uncertainties arising from other causes than source calibration
and counting errors is also briefly discussed, so that the results of this work can be
placed in context with other errors arising in radiometric dosimetry.

\subsection*{A1.2\quad Fitting Functions}

A1.2.1\quad Many papers have been published on the topic of efficiency calibration
functions for germanium spectrometers. Only a few will be mentioned here.

A1.2.2\quad McNelles and Campbell~(9) used a variety of analytical functions, including:

\begin{equation}
  \varepsilon = (a_1/E)^{a_2+a_3}\exp(-a_4 E)
    + a_5\exp(-a_6 E) + a_7\exp(-a_8 E)
  \tag{A1.1}
  \label{eq:A1.1}
\end{equation}

A1.2.3\quad Debertin~(37) introduced the following class of functions:

\begin{equation}
  \varepsilon = a_1\ln(E) + \frac{a_2\ln(E)}{E}
    + \frac{a_3(\ln(E))^2}{E} + \frac{a_4(\ln(E))^4}{E}
    + \frac{a_5(\ln(E))^5}{E}
  \tag{A1.2}
  \label{eq:A1.2}
\end{equation}

A1.2.4\quad Gray and Ahmad~(38) also selected a linear class of fitting function,
defined by:

\begin{equation}
  \varepsilon = \frac{1}{E}\!\left(a_1 + a_2\ln(E) + a_3(\ln(E))^2
    + a_4(\ln(E))^3 + a_5(\ln(E))^5
    + a_6(\ln(E))^7\right)
  \tag{A1.3}
  \label{eq:A1.3}
\end{equation}

A1.2.5\quad Another class of linear fitting functions has been used to fit the logarithm
of the measured efficiencies:

\begin{equation}
  \ln(\varepsilon) = \sum_{i}^{I} a_i \left(\ln\!\left(\frac{E_c}{E}\right)\right)^{\!i-1}
  \tag{A1.4}
  \label{eq:A1.4}
\end{equation}

\noindent where $E_c$ is an arbitrary energy chosen near the middle of the energy range
(to simplify the problem numerically).

A1.2.5.1\quad The number of parameters, $I$, was chosen by Lin et al.~(39) to be six,
for the energy range 300~keV to 9700~keV. Kis et al.~(40) used up to nine parameters for
the range 50~keV to 11\,000~keV. Such a large number is not needed for a narrower energy
range and may be impractical when the number of measured efficiencies is limited. The
number of parameters must always be fewer (ideally many fewer) than the number of
measurements. The purpose of fitting a continuous function to the data is so that
efficiencies at energies in between those included in the calibration can be determined.
Functions that do not sufficiently constrain the values in between the fitted points are
useless for this purpose. Eq~\ref{eq:A1.4}, with $I$ chosen as 6, is used in this
practice.

\subsection*{A1.3\quad Least-Squares Model for Linear Fitting Formulae}

A1.3.1\quad Let $A$ represent a column vector of the parameters $a_i$, for $i = 1, I$.
Then the efficiencies measured for each photon energy of the calibration source can be
represented, using Eq~\ref{eq:A1.4}, by:

\begin{equation}
  \ln(\varepsilon) = S \cdot A + \delta
  \tag{A1.5}
  \label{eq:A1.5}
\end{equation}

\noindent where $\ln(\varepsilon)$ is a column vector of the measured log values for the
efficiencies, $S$ is the design matrix for the model, and $\delta$ is a vector
representing the experimental deviations from the ideal model. When Eq~\ref{eq:A1.4} is
used, the $I$ columns of the matrix $S$ are the values of $(\ln(E_c/E))^{i-1}$,
evaluated at each energy at which the efficiency was measured.

A1.3.2\quad To adapt the linear least-squares model for the case in which
$\log(\text{efficiency})$ is given by a linear model in $A$, we have only to define
$V_{l\varepsilon}$ as the covariance matrix of the logarithms of the measured
efficiencies. If the uncertainties in $\varepsilon$ are assumed to be distributed as log
normal, then the least-squares solution for the parameters $A$ is:

\begin{equation}
  \hat{A} = (S'V_{l\varepsilon}^{-1}S)^{-1}S'V_{l\varepsilon}^{-1}\ln(\varepsilon)
  \tag{A1.6}
  \label{eq:A1.6}
\end{equation}

\noindent with covariance matrix:

\begin{equation}
  V_{\hat{A}} = (S'V_{l\varepsilon}^{-1}S)^{-1}
  \tag{A1.7}
  \label{eq:A1.7}
\end{equation}

A1.3.3\quad The fitted values of the efficiency at the energies of the calibration
measurement are:

\begin{equation}
  \widehat{\ln(\varepsilon)} = S\hat{A}
  \tag{A1.8}
  \label{eq:A1.8}
\end{equation}

\noindent with covariance matrix:

\begin{equation}
  \hat{V}_{l\varepsilon} = S(S'V_{l\varepsilon}^{-1}S)^{-1}S'
  \tag{A1.9}
  \label{eq:A1.9}
\end{equation}

A1.3.4\quad The value of chi-squared must always be calculated when an efficiency fitting
is performed. It is given by:

\begin{equation}
  \chi^2 = \bigl(\ln(\varepsilon) - \widehat{\ln(\varepsilon)}\bigr)'
           \cdot V_{l\varepsilon}^{-1}
           \cdot \bigl(\ln(\varepsilon) - \widehat{\ln(\varepsilon)}\bigr)
  \tag{A1.10}
  \label{eq:A1.10}
\end{equation}

If the number of fitted efficiency values is $N$, then the number of degrees of freedom
of the fit is equal to:

\begin{equation}
  n = N - I
  \tag{A1.11}
  \label{eq:A1.11}
\end{equation}

A1.3.4.1\quad Tables of the chi-squared distribution give the confidence intervals for
the values to be expected assuming that deviations from the nominal value of 1 arise by
chance. For the numerical example given in this annex, $n = 11 - 6 = 5$ and the critical
values for the interval 2.5\,\% to 97.5\,\% of chi-squared for five degrees of freedom
are 0.831 to 12.833. If the observed value is outside this range, then it should be
assumed that the input values and uncertainties are unreliable, and the fit should be
rejected.

A1.3.5\quad Chi-squared values for the test calculated with Eq~\ref{eq:A1.10} will only
be meaningful if the input uncertainties are correctly calculated as one standard
deviation ($1\sigma$). Calibration sources usually have their uncertainties specified
with a confidence of 95\,\% or 99\,\%. According to the GUM~(2), such externally
supplied uncertainties are to be treated as having infinite degrees of freedom, hence a
normal (Gaussian) distribution. Accordingly, the reported calibration uncertainties must
be divided by the $z$-value for a normal distribution that corresponds to the correct
confidence level for a one-tailed test. This must be done before the uncertainties are
combined with $1\sigma$ contributions from other sources, such as counting statistics.
Combination of various uncertainties must be done in quadrature, that is, using the
square root of the sum of squares. A single-lab test of repeatability has shown that the
statistical counting uncertainties may be accurately determined using the Poisson
distribution (Type~B uncertainty estimate). For counts greater than 30, the Poisson
distribution may be approximated with negligible bias by a normal distribution with
standard deviation equal to the square root of the recorded count.

A1.3.6\quad Define $V_{ld}$ as the covariance matrix of the logarithms of the
efficiencies at the gamma energies used to count the activities of activation detectors.
Then:

\begin{equation}
  V_{ld} = D\cdot(S'\cdot V_{l\varepsilon}^{-1}\cdot S)^{-1}\cdot D'
  \tag{A1.12}
  \label{eq:A1.12}
\end{equation}

\noindent where the matrix $D$ is the design matrix that applies for the activation
detector gamma energies. $D$ is formed in the same manner as the $S$ matrix but using
the activation detector gamma energies instead of the calibration source energies.

A1.3.7\quad This result is convenient, because the matrices that specify the log
covariances are the same as the fractional covariances (sometimes called the relative
covariances) of the efficiencies themselves.

\subsection*{A1.4\quad Numerical Example}

A1.4.1\quad A commonly used type of calibration source has eleven gamma energies from
88~keV to 1836~keV, from nine radioisotopes, as shown in Table~\ref{tab:A1.1}.

A1.4.2\quad When Eq~\ref{eq:A1.8} is used to fit measured efficiency values, the
fractional covariance is found using Eq~\ref{eq:A1.9}. The matrix $S$, shown in
Table~\ref{tab:A1.2}, is calculated using the energies in row 2 of Table~\ref{tab:A1.1}.
The matrix $D$ is calculated in the same way using the energies of the activation
detectors. The ones used in the examples here are shown (in~keV) in row 2 of
Table~\ref{tab:A1.3}.

A1.4.3\quad Table~\ref{tab:A1.4} shows the correlation coefficients and percent
uncertainties for the activation detector activities when the Ge spectrometer
efficiencies are fitted with Eq~\ref{eq:A1.4}. The values are calculated for the
activation detector energies of Table~\ref{tab:A1.3}.

A1.4.4\quad The analysis presented shows that activities measured in radiometric monitors
need to have uncertainties no larger than 1.5\,\%, and that those uncertainties are
correlated between detectors because of the detector calibration. Larger uncertainties
are often quoted, and may be present because of filter corrections, flux perturbation
corrections, and positioning errors during irradiation and counting. However, some of
those uncertainties properly belong to the calculations with which the measurements are
compared.

% ---- Table A1.1 ----
\begin{table}[ht]
\centering
\caption{Calibration Source Specifications}
\label{tab:A1.1}
\small
\begin{tabular}{lrrrrrrrrrrr}
\toprule
Radionuclide        & Cd-109 & Co-57 & Ce-139 & Hg-203 & Sn-113 & Sr-85 & Cs-137 & Y-88 & Co-60 & Co-60 & Y-88 \\
Gamma energy (keV)  & 88.03  & 122.1 & 165.9  & 279.2  & 391.7  & 514   & 661.6  & 898  & 1173  & 1333  & 1836 \\
Uncertainty\,\% (1$\sigma$) & 2.9 & 0.8 & 0.8 & 0.8 & 2.1 & 0.8 & 0.9 & 0.8 & 0.8 & 0.8 & 0.8 \\
Correlation coeff.  & \ldots & \ldots & \ldots & \ldots & \ldots & \ldots & \ldots & 0.95+ & 0.999* & 0.999* & 0.95+ \\
\% Counting uncert. & 2.68 & 1.06 & 0.8 & 0.6 & 1.07 & 2.19 & 1.23 & 0.71 & 0.56 & 0.69 & 0.8 \\
\bottomrule
\end{tabular}
\end{table}

% ---- Table A1.2 ----
\begin{table}[ht]
\centering
\caption{The Design Matrix, $S$, for the Fit (elements of $S$ are within Columns 2 to 7)}
\label{tab:A1.2}
\begin{tabular}{ccccccc}
\toprule
$E$ (keV) & $x^0$ & $x^1$ & $x^2$ & $x^3$ & $x^4$ & $x^5$ \\
\midrule
88.03   & 1.0 & 2.3913  & 5.7185  & 13.6748  & 32.7011 & 78.1995  \\
122.06  & 1.0 & 2.0645  & 4.2622  &  8.7993  & 18.1661 & 37.5039  \\
165.85  & 1.0 & 1.7579  & 3.0903  &  5.4326  &  9.5501 & 16.7884  \\
279.19  & 1.0 & 1.2371  & 1.5305  &  1.8934  &  2.3423 &  2.8978  \\
391.69  & 1.0 & 0.8985  & 0.8074  &  0.7255  &  0.6519 &  0.5857  \\
513.99  & 1.0 & 0.6268  & 0.3929  &  0.2463  &  0.1544 &  0.0968  \\
661.65  & 1.0 & 0.3743  & 0.1401  &  0.0524  &  0.0196 &  0.0073  \\
898.02  & 1.0 & 0.0688  & 0.0047  &  0.0003  &  0.0000 &  0.0000  \\
1173.22 & 1.0 & $-$0.1985 & 0.0394 & $-$0.0078 & 0.0016 & $-$0.0003 \\
1332.49 & 1.0 & $-$0.3258 & 0.1061 & $-$0.0346 & 0.0113 & $-$0.0037 \\
1836.01 & 1.0 & $-$0.6463 & 0.4177 & $-$0.2700 & 0.1745 & $-$0.1128 \\
\bottomrule
\end{tabular}
\end{table}

% ---- Table A1.3 ----
\begin{table}[ht]
\centering
\caption{Activation Detector Activities and Energies in keV}
\label{tab:A1.3}
\small
\begin{tabular}{llllllllllllllllll}
\toprule
Sc47 & In115m & Au198 & Ag110m & W187 & Co58 & Mn54 & Mn56 & Sc46 & Zr89 & Nb92m & Sc48 & Fe59 & Co60 & In116 & Cu64 & Na24 & Ni57 \\
159 & 336 & 412 & 658 & 686 & 811 & 835 & 847 & 889 & 909 & 935 & 984 & 1099 & 1173 & 1294 & 1346 & 1369 & 1378 \\
\bottomrule
\end{tabular}
\end{table}

% ---- Table A1.4 ----
\begin{landscape}
\begin{table}[p]
\centering
\caption{Percentage Uncertainties and Correlation Coefficients for Measured Activities of
Isotopes Shown in Table~\ref{tab:A1.3}, Using Photons at the Energies Listed in keV (row 2 of
Table~\ref{tab:A1.3})}
\label{tab:A1.4}
\begin{astmnote}[1]
Uncertainties attributable to calibration source strength and counting errors are as shown
in Table~\ref{tab:A1.1}. Counting errors of 1\,\% were assumed on each of the peak areas
measured in activity measurements.
\end{astmnote}
\small
\setlength{\tabcolsep}{3pt}
\begin{tabular}{rr rrrrrrrrrrrrrrrrrr}
\toprule
s.d. & keV & 159 & 336 & 412 & 658 & 686 & 811 & 835 & 847 & 889 & 909 & 935 & 984 & 1099 & 1173 & 1294 & 1346 & 1369 & 1378 \\
\midrule
1.33\,\% & 159  & 1.00  & $-$0.01 & $-$0.07 & $-$0.05 & $-$0.05 & $-$0.02 & $-$0.01 & $-$0.01 & 0.00 & 0.00 & 0.00 & 0.01 & 0.02 & 0.02 & 0.02 & 0.01 & 0.01 & 0.01 \\
1.32\,\% & 336  & $-$0.01 & 1.00 & 0.40 & 0.16 & 0.14 & 0.06 & 0.05 & 0.04 & 0.02 & 0.01 & 0.00 & $-$0.02 & $-$0.04 & $-$0.04 & $-$0.04 & $-$0.04 & $-$0.04 & $-$0.04 \\
1.36\,\% & 412  & $-$0.07 & 0.40 & 1.00 & 0.34 & 0.32 & 0.21 & 0.18 & 0.17 & 0.13 & 0.11 & 0.09 & 0.04 & $-$0.05 & $-$0.09 & $-$0.12 & $-$0.13 & $-$0.13 & $-$0.13 \\
1.40\,\% & 658  & $-$0.05 & 0.16 & 0.34 & 1.00 & 0.48 & 0.41 & 0.38 & 0.37 & 0.32 & 0.30 & 0.27 & 0.20 & 0.05 & $-$0.03 & $-$0.11 & $-$0.14 & $-$0.14 & $-$0.14 \\
1.38\,\% & 686  & $-$0.05 & 0.14 & 0.32 & 0.48 & 1.00 & 0.41 & 0.39 & 0.38 & 0.33 & 0.31 & 0.28 & 0.22 & 0.07 & $-$0.01 & $-$0.09 & $-$0.12 & $-$0.12 & $-$0.13 \\
1.28\,\% & 811  & $-$0.02 & 0.06 & 0.21 & 0.41 & 0.41 & 1.00 & 0.38 & 0.38 & 0.35 & 0.34 & 0.32 & 0.28 & 0.17 & 0.11 & 0.03 & 0.01 & 0.00 & 0.00 \\
1.26\,\% & 835  & $-$0.01 & 0.05 & 0.18 & 0.38 & 0.39 & 0.38 & 1.00 & 0.37 & 0.35 & 0.34 & 0.32 & 0.28 & 0.19 & 0.13 & 0.06 & 0.04 & 0.03 & 0.03 \\
1.26\,\% & 847  & $-$0.01 & 0.04 & 0.17 & 0.37 & 0.38 & 0.38 & 0.37 & 1.00 & 0.35 & 0.34 & 0.32 & 0.29 & 0.20 & 0.14 & 0.07 & 0.05 & 0.04 & 0.04 \\
1.23\,\% & 889  & 0.00 & 0.02 & 0.13 & 0.32 & 0.33 & 0.35 & 0.35 & 0.35 & 1.00 & 0.33 & 0.32 & 0.30 & 0.23 & 0.19 & 0.13 & 0.11 & 0.10 & 0.10 \\
1.22\,\% & 909  & 0.00 & 0.01 & 0.11 & 0.30 & 0.31 & 0.34 & 0.34 & 0.34 & 0.33 & 1.00 & 0.32 & 0.31 & 0.25 & 0.21 & 0.15 & 0.13 & 0.13 & 0.12 \\
1.21\,\% & 935  & 0.00 & 0.00 & 0.09 & 0.27 & 0.28 & 0.32 & 0.32 & 0.32 & 0.32 & 0.32 & 1.00 & 0.31 & 0.27 & 0.23 & 0.19 & 0.17 & 0.16 & 0.16 \\
1.21\,\% & 984  & 0.01 & $-$0.02 & 0.04 & 0.20 & 0.22 & 0.28 & 0.28 & 0.29 & 0.30 & 0.31 & 0.31 & 1.00 & 0.30 & 0.28 & 0.25 & 0.23 & 0.23 & 0.22 \\
1.24\,\% & 1099 & 0.02 & $-$0.04 & $-$0.05 & 0.05 & 0.07 & 0.17 & 0.19 & 0.20 & 0.23 & 0.25 & 0.27 & 0.30 & 1.00 & 0.36 & 0.36 & 0.35 & 0.35 & 0.35 \\
1.28\,\% & 1173 & 0.02 & $-$0.04 & $-$0.09 & $-$0.03 & $-$0.01 & 0.11 & 0.13 & 0.14 & 0.19 & 0.21 & 0.23 & 0.28 & 0.36 & 1.00 & 0.41 & 0.41 & 0.40 & 0.40 \\
1.34\,\% & 1294 & 0.02 & $-$0.04 & $-$0.12 & $-$0.11 & $-$0.09 & 0.03 & 0.06 & 0.07 & 0.13 & 0.15 & 0.19 & 0.25 & 0.36 & 0.41 & 1.00 & 0.45 & 0.45 & 0.45 \\
1.36\,\% & 1346 & 0.01 & $-$0.04 & $-$0.13 & $-$0.14 & $-$0.12 & 0.01 & 0.04 & 0.05 & 0.11 & 0.13 & 0.17 & 0.23 & 0.35 & 0.41 & 0.45 & 1.00 & 0.46 & 0.46 \\
1.37\,\% & 1369 & 0.01 & $-$0.04 & $-$0.13 & $-$0.14 & $-$0.12 & 0.00 & 0.03 & 0.04 & 0.10 & 0.13 & 0.16 & 0.23 & 0.35 & 0.40 & 0.45 & 0.46 & 1.00 & 0.46 \\
1.37\,\% & 1378 & 0.01 & $-$0.04 & $-$0.13 & $-$0.14 & $-$0.13 & 0.00 & 0.03 & 0.04 & 0.10 & 0.12 & 0.16 & 0.22 & 0.35 & 0.40 & 0.45 & 0.46 & 0.46 & 1.00 \\
\bottomrule
\end{tabular}
\end{table}
\end{landscape}

% ============================================================
\section*{References}

\begin{enumerate}[label=(\arabic*),leftmargin=3em]
  \item Taylor, B. and Kuyatt, C., ``Guidelines for Evaluating and Expressing the
        Uncertainty of NIST Measurement Results,'' NIST Technical Note 1297, 1994.

  \item ``Evaluation of measurement data---Guide to the expression of uncertainty in
        measurement,'' JCGM 100:2008, GUM 1995 with minor corrections, BIPM, 2008.

  \item Helmer, R.~G., et al., \textit{Nuclear Instruments and Methods}, Vol~96, 1972,
        p.~173.

  \item Evaluated Nuclear Structure Data File (ENSDF), a computer file of evaluated
        nuclear structure and radioactive decay data, which is maintained by the National
        Nuclear Data Center (NNDC), Brookhaven National Laboratory (BNL), on behalf of
        The International Network for Nuclear Structure Data Evaluation, which functions
        under the auspices of the Nuclear Data Section of the International Atomic Energy
        Agency (IAEA).

  \item \textit{Nuclear Wallet Cards}, compiled by J.~K. Tuli, National Nuclear Data
        Center, \url{www.nndc.bnl.gov/wallet/wccurrent.html}.

  \item Freeman, J.~M., and Jenkin, J.~G., \textit{Nuclear Instruments and Methods},
        Vol~43, 1966, p.~269.

  \item Tokcan, G., and Cothern, C.~R., \textit{Nuclear Instruments and Methods},
        Vol~64, 1968, p.~219.

  \item Paradellis, T., and Hontzeas, S., \textit{Nuclear Instruments and Methods},
        Vol~73, 1969, p.~210.

  \item McNelles, L.~A., and Campbell, J.~L., \textit{Nuclear Instruments and Methods},
        Vol~109, 1973, p.~241.

  \item Debertin, K., et al., \textit{Proceedings of ERDA Symposium on X- and Gamma-Ray
        Sources and Applications}, CONF-760539, 1976, p.~59. Available from the National
        Technical Information Service.

  \item Hirschfeld, A.~T., et al., \textit{Proceedings of ERDA Symposium on X- and
        Gamma-Ray Sources and Applications}, CONF-760539, 1976, p.~90. Available from
        the National Technical Information Service.

  \item Singh, R., \textit{Nuclear Instruments and Methods}, Vol~136, 1976, p.~543.

  \item Kawade, K., et al., \textit{Nuclear Instruments and Methods}, Vol~190, 1981,
        p.~101.

  \item Moens, L., et al., \textit{Nuclear Instruments and Methods}, Vol~187, 1981,
        p.~451.

  \item Koskelo, M.~J., \textit{Journal of Radioanalytical Chemistry}, Vol~70, 1982,
        p.~473.

  \item Houtermans, H., et al., \textit{International Journal of Applied Radiation and
        Isotopes}, Vol~34, 1983, p.~487.

  \item Wyttenbach, A., \textit{Journal of Radioanalytical Chemistry}, Vol~8, 1971,
        p.~335.

  \item Cohen, E.~J., \textit{Nuclear Instruments and Methods}, Vol~121, 1974, p.~25.

  \item Strauss, M.~G., et al., \textit{IEEE Transactions of Nuclear Science NS-15},
        1968, p.~518.

  \item Anders, O.~U., \textit{Nuclear Instruments and Methods}, Vol~68, 1969, p.~205.

  \item Wiernik, M., \textit{Nuclear Instruments and Methods}, Vol~96, 1971, p.~325.

  \item Debertin, K., and Schotzig, U., \textit{Nuclear Instruments and Methods},
        Vol~140, 1977, p.~337.

  \item Westphal, G.~P., \textit{Nuclear Instruments and Methods}, Vol~163, 1979,
        p.~189.

  \item Routti, J.~T., and Prussin, S.~G., \textit{Nuclear Instruments and Methods},
        Vol~72, 1969, p.~125.

  \item Gunnink, R., and Niday, J.~B., University of California, Lawrence Livermore
        Laboratory Report, UCRL-51061, 1972.

  \item Yule, H.~P., \textit{Analytical Chemistry}, Vol~40, 1968, p.~1480.

  \item Heydorn, K., and Lada, W., \textit{Analytical Chemistry}, Vol~44, 1972, p.~2313.

  \item Kokta, L., \textit{Nuclear Instruments and Methods}, Vol~112, 1973, p.~245.

  \item Helmer, R.~G., and Lee, M.~A., \textit{Nuclear Instruments and Methods},
        Vol~178, 1980, p.~499.

  \item Schima, F.~J., and Hoppes, D.~D., \textit{International Journal of Applied
        Radiation and Isotopes}, Vol~34, 1983, p.~1109.

  \item Debertin, K., and Schotzig, U., \textit{Nuclear Instruments and Methods},
        Vol~158, 1979, p.~471.

  \item Morel, J., et al., \textit{International Journal of Applied Radiation and
        Isotopes}, Vol~34, 1983, p.~1115.

  \item McCallum, G.~J., and Coote, G.~E., \textit{Nuclear Instruments and Methods},
        Vol~130, 1975, p.~189.

  \item Hubbell, J.~H., and Seltzer, S.~M., ``X-Ray Mass Attenuation Coefficients,''
        NIST Standard Reference Database 126,
        \url{www.nist.gov/pml/x-ray-mass-attenuation-coefficients}.

  \item Williams, J.~G., ``Covariances for Gamma Spectrometer Efficiency Calibrations,''
        \textit{Fifteenth International Symposium on Reactor Dosimetry},
        Aix-en-Provence, 2014, EPJ Web of Conferences, 106, 07004 (2016).

  \item Seber, G.~A.~F., \textit{Linear Regression Analysis}, Wiley, New York, 1977.

  \item Debertin, K., ``The Effect of Correlations in the Efficiency Calibration of
        Germanium Detectors,'' \textit{Nuclear Instruments and Methods}, Vol~226, 1984,
        pp.~566--568.

  \item Gray, P.~W., and Ahmad, A., ``Linear Classes of Ge(Li) Detector Efficiency
        Functions,'' \textit{Nuclear Instruments and Methods A}, Vol~237, 1985,
        pp.~577--589.

  \item Lin, J., Henry, E., and Meyer, R.~A., Lawrence Livermore National Laboratory
        Report UCLR-84536, 1980.

  \item Kis, Z., et al., ``Comparison of efficiency functions for Ge gamma-ray detectors
        in a wide energy range,'' \textit{Nuclear Instruments and Methods in Physics
        Research A}, Vol~418, 1998, pp.~374--386.
\end{enumerate}

% ============================================================
\vspace{2em}
\noindent\rule{\textwidth}{0.4pt}

\noindent\footnotesize
$^{1}$ This practice is under the jurisdiction of ASTM Committee E10 on Nuclear
Technology and Applications and is the direct responsibility of Subcommittee E10.05 on
Nuclear Radiation Metrology. Current edition approved Feb.~15, 2023. Published April
2023. DOI: 10.1520/E3376-23.

\noindent
$^{2}$ The boldface numbers in parentheses refer to a list of references at the end of
this standard.

\noindent
$^{3}$ For referenced ASTM standards, visit the ASTM website, \url{www.astm.org}, or
contact ASTM Customer Service at \href{mailto:service@astm.org}{service@astm.org}.

\noindent
$^{4}$ Available from Institute of Electrical and Electronics Engineers, Inc. (IEEE),
445 Hoes Ln., Piscataway, NJ 08854-4141, \url{http://www.ieee.org}.

\noindent
$^{5}$ For simplicity of these calculations, $n$ is assumed to be the same on both sides
of the peak. If the continuum is calculated using a different number of channels on the
left of the peak than on the right of the peak, different equations must be used.

\end{document}
